\begin{problem}{Подземелье для принцесс}{стандартный ввод}{стандартный вывод}{2 секунды}{256 мегабайт}

Боузер~--- злодей, который часто похищает принцесс, и даже имеет для их заточения целое подземелье. 

В подземелье расположено $n$ тюремных камер. Камеры пронумерованы натуральными числами от $1$ до $n$ и расположены в этом порядке вдоль длинного коридора на одинаковом расстоянии друг от друга. Камера с номером $i$ рассчитана на заточение $a_i$ принцесс. Вход в подземелье расположен между камерами $k$ и $k + 1$, при этом, если $k = 0$, это означает, что вход находится перед первой камерой. Если $k = n$, вход расположен после последней камеры.

Боузер собирается сделать $m$ вылазок с целью похищения принцесс. В $j$-ю вылазку он планирует похитить $b_j$ принцесс. Приведя их в подземелье, 
Боузер выбирает камеру для новых принцесс по следующему принципу:
\begin{itemize}
\item Если есть свободная камера с вместимостью $a_i = b_j$, то он обязательно выбирает такую камеру. Иначе он выбирает свободную камеру с вместимостью $a_i > b_j$.
\item Из подходящих камер он выбирает ту, которая расположена как можно ближе ко входу. Если камер с одинаковым расстоянием до входа несколько, он выбирает ту, у которой номер меньше. Расстояние от входа до камеры Боузер считает равным количеству камер между ними.
\item Если свободной камеры, подходящей по вместимости, нет, то ему остается лишь расстроиться и отпустить принцесс.
\end{itemize}
Чтобы не допустить побега принцесс, после того как Боузер заточил их в камеру, он закрывает камеру на ключ и выбрасывает его. Таким образом открыть камеру и добавить туда принцесс похищенных в другой вылазке, невозможно. Также Боузер никогда не пытается разместить принцесс, похищенных в одной вылазке, более чем в одной камере.

Марио и Луиджи нашли планы подземелья и владеют информацией о вылазках Боузера. Для каждой вылазки, они хотят определить камеру, в которую будут заточены принцессы.

\InputFile
В первой строке даны три целых числа $n$, $m$ и $k$~--- число тюремных камер, число вылазок Боузера и номер камеры, после которой расположен вход в подземелье ($1 \le n, m \le 1000$, $0 \le k \le n$). Если $k = 0$, то вход находится перед первой камерой.

Во второй строке даны $n$ целых чисел $a_i$~--- вместительности камер ($1 \le a_i \le 1000$).

Во третьей строке даны $m$ целых чисел $b_j$~--- число похищенных в $j$-ю вылазку принцесс ($1 \le b_j \le 1000$).

\OutputFile
Выведите $m$ чисел $c_j$~--- номера тюремных камер, выбранных для принцесс, похищенных в $j$-ю вылазку. Если Боузер отпустил принцесс, так как не нашел для них подходящую камеру, будем считать, что $c_j = -1$.

\Scoring
Эта задача состоит из четырех подзадач. Для подзадач выполняются дополнительные ограничения, указанные в таблице ниже. Для получения баллов за подзадачу необходимо пройти все тесты данной подзадачи, а также все тесты всех необходимых подзадач. Номера необходимых подзадач также указаны в таблице.

\begin{center}
\renewcommand{\arraystretch}{1.5}
\begin{tabular}{|c|c|p{0.4\textwidth}|c|}
\hline
Подзадача & Баллы & Дополнительные ограничения & Необходимые подзадачи
\\ \hline
1 & 21 & $k = 0$ & {---}
\\ \hline

2 & 22 & $a_i = 1$ & {---}
\\ \hline

3 & 23 & $b_j = 1$ & {---}
\\ \hline
4 & 34 & Полные ограничения & 1, 2, 3
\\ \hline

\end{tabular}

\end{center}

\Examples

\begin{example}
\exmp{3 3 0
5 3 1
1 1 1
}{3 1 2 
}%
\exmp{4 5 2
1 3 2 2
1 1 3 2 1
}{1 2 -1 3 4 
}%
\exmp{6 6 3
1 2 3 3 2 1
1 1 1 1 1 1
}{1 6 3 4 2 5 
}%
\end{example}

\end{problem}

